\chapter{Abdominal Aortic Aneurysm}
\index{Abdominal Aortic Aneurysm}
\label{sec:Abdominal Aortic Aneurysm}

\section{Overview}
An abdominal aortic aneurysm (AAA) is a permanent enlargement of the abdominal aorta, usually defined as an aortic diameter of at least 3 cm or more than 50\% above the expected normal size. The aneurysm may enlarge silently over years. The principal danger is rupture, which causes life-threatening internal hemorrhage. Management is based on aneurysm size, growth rate, symptoms, operative risk, and anatomy.
\section{Causes}
AAA results from progressive weakening and remodeling of the aortic wall. Atherosclerotic inflammation, loss of elastic tissue, and degradation of structural proteins contribute to wall dilation. Less common causes include inherited connective-tissue disorders, vasculitis, infection, trauma, and prior aortic surgery.
\section{Risk Factors}
\begin{itemize}[noitemsep]
  \item Increasing age, especially older age in men
  \item Current or previous cigarette smoking
  \item Family history of AAA
  \item Hypertension and atherosclerotic cardiovascular disease
  \item Male sex and other vascular aneurysms
  \item Genetic connective-tissue disorders
\end{itemize}
\section{Signs and Symptoms}
Most AAAs are asymptomatic and are discovered during examination or imaging for another reason. Possible symptoms include persistent or deep abdominal, flank, or back pain and a pulsating abdominal mass. Sudden severe abdominal or back pain, fainting, low blood pressure, pallor, or collapse may indicate rupture and require emergency treatment.
\section{Progression and Stages}
Small aneurysms may remain stable and can be monitored with periodic ultrasound. Some enlarge gradually, while others expand more rapidly. Symptoms, rapid growth, or a diameter at which rupture risk outweighs procedural risk generally prompt repair. Rupture can lead to shock and death within a short time.
\section{Complications}
\begin{itemize}[noitemsep]
  \item Aneurysm rupture and internal bleeding
  \item Thrombus formation with embolization to the legs or organs
  \item Compression of nearby abdominal structures
  \item Aortoenteric or aortocaval fistula, rarely
  \item Kidney, bowel, or limb ischemia
\end{itemize}
\section{Diagnosis}
Abdominal ultrasound is commonly used for detection and surveillance. Computed tomography angiography defines the aneurysm's size, extent, branch-vessel involvement, and suitability for repair. Magnetic resonance angiography may be used when CT contrast or radiation is unsuitable. Evaluation also includes cardiovascular risk assessment and preoperative testing.
\section{Prevention}
\begin{itemize}[noitemsep]
  \item Stop smoking and avoid second-hand smoke
  \item Control blood pressure, diabetes, and lipid abnormalities
  \item Maintain regular physical activity and a healthy weight
  \item Attend scheduled imaging surveillance
  \item Seek urgent care for sudden severe abdominal or back pain
\end{itemize}
\section{Treatment Options}
Small, asymptomatic AAAs are usually managed with surveillance and risk-factor reduction. Symptomatic, rapidly enlarging, or sufficiently large aneurysms require vascular-surgery assessment for elective repair. Endovascular aneurysm repair places a stent graft inside the aorta, whereas open repair replaces the diseased segment with a graft. A suspected rupture requires immediate resuscitation and emergency surgical or endovascular repair.
\section{Living with It}
Patients should keep imaging and vascular-surgery appointments, take prescribed cardiovascular medicines, stop smoking, and report new abdominal, flank, or back pain promptly. Physical activity is generally encouraged, but the treating clinician should guide activity when the aneurysm is large or symptomatic.
\section{Outlook}
The outlook is favorable for many patients when an AAA is detected before rupture and monitored or repaired at an appropriate time. Prognosis after rupture is substantially worse, making surveillance and rapid attention to warning symptoms essential.
